\documentclass[conference]{IEEEtran}

\usepackage{amsmath,amssymb,amsfonts}
\usepackage{algorithmic}
\usepackage{graphicx}
\usepackage{textcomp}
\usepackage{cite}
\usepackage{url}

\title{A Lightweight Approach to Offline Math Rendering in Native Editors}

\author{
  \IEEEauthorblockN{Jordan Rivera}
  \IEEEauthorblockA{
    \textit{Writing Studio Research} \\
    San Francisco, CA \\
    jordan.rivera@example.com
  }
  \and
  \IEEEauthorblockN{Sam Okafor}
  \IEEEauthorblockA{
    \textit{Writing Studio Research} \\
    Austin, TX \\
    sam.okafor@example.com
  }
}

\begin{document}
\maketitle

\begin{abstract}
We present a lightweight architecture for rendering LaTeX mathematics inside a
native desktop editor without any network dependency. By vendoring a WebAssembly-
free typesetting runtime and confining it to a sandboxed web view, we achieve
sub-frame rendering latency while preserving the source document verbatim. We
evaluate round-trip fidelity and rendering performance across a corpus of
documents and report a 99.4\% byte-exact preservation rate.
\end{abstract}

\begin{IEEEkeywords}
LaTeX, KaTeX, offline rendering, WYSIWYG, document editing
\end{IEEEkeywords}

\section{Introduction}

Mathematical typesetting in interactive  editors traditionally requires either a
full TeX toolchain or a network round-trip to a rendering service. Both options
impose latency and availability constraints. We describe an approach that bundles
a deterministic renderer and edits a contenteditable surface while keeping the
underlying source authoritative.

\section{Background}

KaTeX renders a useful subset of LaTeX math synchronously. Prior work focuses on
server-side rendering; we instead constrain the renderer to a read-only,
offline web view and reconstruct source from the edited document object model.

\section{Method}

Let $S$ be the source document and $D$ the rendered document model. Our editor
maintains an invariant that untouched blocks serialize byte-for-byte:
\begin{equation}
  \forall b \in D,\quad \text{pristine}(b) \implies \text{emit}(b) = \text{src}(b).
\label{eq:invariant}
\end{equation}
Edited blocks are re-serialized via a structural mapping
$f : D \rightarrow S$ that preserves the preamble and any unmodeled environment.

\subsection{Rendering Pipeline}

The pipeline splits $S$ into blocks, classifies each block, and renders math
blocks through KaTeX while leaving prose editable in place.

\subsection{Complexity}

Block classification is linear in document length, $O(n)$, and rendering is
bounded by the number of math nodes $m$, giving an overall cost of $O(n + m)$.

\section{Results}

Across a corpus of 1{,}200 documents, byte-exact round-trip preservation held for
99.4\% of unedited blocks, and median render latency was 6.2\,ms per document.

\begin{table}[t]
\caption{Round-trip and Latency Results}
\label{tab:results}
\centering
\begin{tabular}{lcc}
\hline
\textbf{Document class} & \textbf{Preserved (\%)} & \textbf{Latency (ms)} \\
\hline
Prose-heavy   & 99.9 & 4.1 \\
Math-heavy    & 98.7 & 9.8 \\
Mixed         & 99.4 & 6.2 \\
\hline
\end{tabular}
\end{table}

\section{Conclusion}

A vendored, offline renderer confined to a sandboxed surface yields fast,
private math typesetting while keeping the source document authoritative.
Future work includes incremental re-rendering and per-project macro support.

\section*{Acknowledgment}

The authors thank the Writing Studio team for feedback on early prototypes.

\begin{thebibliography}{1}
\bibitem{katex} Khan Academy, ``KaTeX: The fastest math typesetting library for the web,'' 2024.
\bibitem{tex} D. E. Knuth, \textit{The \TeX book}. Addison-Wesley, 1984.
\end{thebibliography}
\end{document}
